\section{Open questions from this replication}

Five open questions raised (or left unresolved) by this replication effort,
each with a concrete, free-endpoint next step.

\subsection*{Q1. Does UCCSD-VQE $\approx$ CCSD survive for the paper's open-shell / strongly-correlated cases?}
\textbf{Basis.} Paper reports UCCSD-VQE $\approx$ CCSD for all 9 molecules
(closed- and open-shell); this replication verified only closed-shell
H$_2$ and LiH. For open-shell :CH$_2$ / CH$_2$ singlet-triplet gap and
stretched N$_2$, single-reference CCSD is documented to break down, so
the analytical identity UCCSD-VQE $=$ CCSD would inherit that failure.

\textbf{Next step.} On qiskit-nature $+$ PySCF, run classical UCCSD (proxy
for converged UCCSD-VQE) and FCI for :CH$_2$ singlet, CH$_2$ triplet,
N$_2$ at $r=1.1, 1.5, 2.0$ \AA{} in STO-3G; tabulate
$|E_\text{UCCSD}-E_\text{FCI}|$ in kJ/mol; flag any case exceeding
$4$ kJ/mol chemical accuracy. Free endpoint (local CPU).

\subsection*{Q2. Does modern low-rank / tensor-hypercontraction (THC) change the paper's CNOT-count estimate?}
\textbf{Basis.} Paper counts CNOTs with MP2 pre-screening $+$ gate
cancellation. Post-2020 (von Burg et al.\ 2021, Lee/Berry et al.\ 2021 THC
for QPE, Motta et al.\ 2021 DF/CDF for VQE) has shown 10--100$\times$
reductions in gate count from integral factorization alone. Paper's
near-term-feasibility conclusion (``10$^5$--10$^7$ CNOTs, out of reach'')
was pre-THC.

\textbf{Next step.} Reimplement LiH and H$_2$O STO-3G Hamiltonian in
OpenFermion $+$ PennyLane with DF/CDF factorization; compare transpiled
CNOT counts to paper Table SI\,I. If reductions $\ge 3\times$ hold, note
paper's feasibility conclusion may need revision. Free endpoint.

\subsection*{Q3. Would active-space UCCSD-VQE dramatically shrink the paper's qubit budget?}
\textbf{Basis.} Paper uses full-orbital JW (6 spatial $\to$ 12 qubits for
LiH STO-3G). Modern practice (ADAPT-VQE, iQCC) uses active-space
partitioning to fit real molecules in $<10$ qubits without losing
chemical accuracy. Paper does not discuss this trade-off; its resource
Pareto may be pessimistic by an order of magnitude in qubit count.

\textbf{Next step.} Run CAS(4,4) $+$ UCCSD-VQE for LiH and H$_2$O STO-3G
on qiskit-nature; report qubit-count vs $|E-E_\text{FCI}|$ Pareto curve
against full-space numbers from this replication $+$ paper. Free endpoint.

\subsection*{Q4. How does UCCSD-VQE stack up against classical near-exact methods (DMRG, CCSD(T), CCSDT)?}
\textbf{Basis.} Paper's implicit pitch is ``VQE gives FCI-quality
answers.'' But for the 9 small molecules in Table SI\,I, DMRG with modest
bond dimension reaches FCI to sub-mHa at laptop cost; CCSD(T) reaches
chemical accuracy for closed-shell cases. Paper does not benchmark
against these, so the quantum-advantage argument is left implicit and
unquantified.

\textbf{Next step.} Using pyscf $+$ block2 (or CheMPS2), run
DMRG($M=500,1000,2000$) for LiH, H$_2$O, N$_2$ STO-3G; compare wall-clock
and $|E-E_\text{FCI}|$ against this replication's UCCSD-VQE numbers. If
DMRG hits $<0.1$ mHa in $<1$ s on a laptop, that reframes the resource
story (compare against DMRG, not FCI). Free endpoint.

\subsection*{Q5. Under a fixed total-shot budget, does measurement-optimized VQE actually beat QPE?}
\textbf{Basis.} Paper focuses on gate-count resources for UCCSD-VQE
circuits without treating the shot cost of VQE energy evaluation
(10$^6$--10$^9$ shots per energy at chemical accuracy). QPE has
poly-log shot cost but higher per-run gate cost. Paper's
near-term-feasibility conclusion depends on which trade-off wins.

\textbf{Next step.} For LiH STO-3G, estimate (a) UCCSD-VQE shots needed
for $E$ to 1 mHa via grouped-Pauli measurement (OpenFermion utilities),
(b) QPE ancilla $+$ T-count for the same target via Berry et al.\ 2019
trotterized-block-encoding recipes; report crossover total-run-time vs
physical-gate error rate; flag which end of the trade-off the paper
implicitly sits on. Free endpoint (OpenFermion, local).
